\section{Экономика}
\label{sec:economics}

Все числа этого раздела посчитаны движком расчётов из репозитория статьи. Каждое входное значение там помечено как
замер, как цена с источником или как оценка. Профиль нагрузки тяжёлый: 2\,500 вызовов в день на человека, 21 рабочий
день в месяц, восьмичасовой рабочий день. На открытую модель уходит доля $s$ этих вызовов (раздел~\ref{sec:workload}).
Модель Qwen3.8-27B, сервер считает в тех же 4-битных весах, что и клиент. Цены API взяты у поставщиков этой же модели: на нашем профиле
1000 вызовов у медианного поставщика стоят \$9.68, у самого дешёвого \$1.34. Генерация на вызов равна ответу по
журналам, 512 токенов, в которые рассуждения модели уже входят (раздел~\ref{sec:workload}), а прошлый ответ модели
оплачивается как новый ввод. Если прибавить рассуждения к ответу ещё раз, получается верхняя граница генерации:
ёмкость M5 Max падает с 883 до 637 вызовов в день, API дорожает на 5--20\%, соотношение с медианным поставщиком
меняется мало (при доле 50\% гибрид \$49 против \$268), а разрыв с самым дешёвым сокращается: \$49 против \$42 при
доле 50\% и в 1.5 раза при доле 20\%.

\textbf{Годовой счёт.} Для масштаба мы оценили по ценам API октября 2026 года сохранившиеся журналы
агентских сессий за год: 1.53 миллиона вызовов агентов, у Codex с ноября 2025 года, у Claude с мая 2026 года, у Muse с
августа 2026 года (таблица~\ref{tab:owner}). По спискам цен они стоили бы \$113--124 тысячи. Вызов Claude обходится
в среднем в 16 центов, вызов Codex в 6--7 центов; вызов в гибриде с Qwen3.8-27B стоит 0.23--0.33 цента. По индексу
Artificial Analysis на моделях уровня Qwen3.8-27B и ниже прошло 1.4\% вызовов, ещё 15.6\% на моделях, которые не
выше Qwen3.8-27B на стандартном усилии рассуждения. Вместе это не больше 17\% вызовов (с каким усилием они шли, журналы не
фиксируют), близко к доле $s = 20\%$ наших расчётов, но только 5\% счёта: такие вызовы шли на дешёвых моделях. Основные деньги приходятся на передовые модели.
Классы здесь определены по модели; какие задачи передовых моделей можно отдать открытой модели, по журналам сказать
нельзя. Ещё 6.2 тысячи вызовов в журналах уже выполнены на открытых моделях локально на Маке.

\begin{table}[t]
\centering
\caption{Сохранившиеся журналы агентских сессий за год по ценам API октября 2026 года, по классам моделей. Класс задан
индексом Artificial Analysis на той же версии индекса, что и у Qwen3.8-27B. Диапазоны отражают изменения цен за год
и неизвестные в журналах величины, например смысл кэша у Muse. Запись в кэш у Codex журналы показывают нулевой на
каждом вызове, и она так и посчитана. Данные и расчёт: \texttt{data/owner-model-usage-2026-10-30.json},
\texttt{data/year\_bill.py}.}
\label{tab:owner}
\footnotesize
\begin{tabular}{lrrr}
\toprule
Класс моделей & Вызовов & Доля & Счёт \\
\midrule
Не выше Qwen3.8-27B & 21\,935 & 1.4\% & \$0.6 тыс. \\
Не выше на стандартном усилии рассуждения & 238\,266 & 15.6\% & \$4.8--5.1 тыс. \\
Выше Qwen3.8-27B & 1\,248\,111 & 81.5\% & \$107--118 тыс. \\
Без индекса & 13\,157 & 0.9\% & \$0.3 тыс. \\
Открытые модели локально на Маке & 6\,178 & 0.4\% & --- \\
Без цены в списке & 3\,092 & 0.2\% & --- \\
\midrule
Всего & 1\,530\,739 & 100\% & \$113--124 тыс. \\
\bottomrule
\end{tabular}
\end{table}

\textbf{Мак уже есть.} Пусть у пользователя уже есть машина класса Max. Тогда гибрид добавляет к его расходам
серверный prefill, хранение копий сессий на сервере, трафик, электричество клиента и обслуживание сервиса. Серверный
prefill мы оплачиваем по часовой цене B300, \$9.50 в час при загрузке 70\%, трафик по \$5 за терабайт, обслуживание
оцениваем в \$5 на человека в месяц. Хранение считаем по политике, которую допускает протокол: сервер дописывает в
копию сессии только новые позиции, держит копии в оперативной памяти в рабочие часы и на NVMe в остальное время, а токены вызовов, прочитанных клиентом, дочитывает при следующем оффлоаде. Гигабайт в месяц мы оцениваем в \$0.50 в оперативной
памяти и \$0.10 на NVMe. Для M5 Max при доле 20\% хранение стоит \$6.5 в месяц, больше, чем сам
prefill (\$5.1). Осторожная политика, при которой копии весь месяц лежат в оперативной памяти, а сессия на каждом
вызове перекачивается на GPU, обходится в \$16.5. Таблица~\ref{tab:econ} показывает итог для разных клиентов.

При доле 20\% пользователь отправляет на открытую модель 500 вызовов в день, и клиенты класса Max почти все их
обслуживают сами. Гибрид на клиентах классов Max и Ultra обходится в \$22--27 в месяц. Это в 1.6--1.9 раза дороже, чем у самого дешёвого поставщика API
(\$14), и в 3.7--4.6 раза дешевле, чем у медианного (\$102). При осторожной политике
хранения гибрид в 2.3--2.6 раза дороже самого дешёвого поставщика и в 2.8--3.1 раза дешевле медианного.

При доле 50\% на открытую модель уходит 1\,250 вызовов в день, и здесь проявляется предел клиента. Генерация
занимает время, и за восьмичасовой день M5 Max с четырьмя потоками успевает обслужить около 880 вызовов. Остальные мы
отправляем самому дешёвому поставщику API и включаем их стоимость в цену гибрида. Выходит \$43--49 в месяц: на 23--38\% дороже самого
дешёвого API (\$35) и в 5.2--5.9 раза дешевле медианного (\$254). Это соотношение складывает два
эффекта: гибрид и переход к самому дешёвому поставщику для вызовов сверх ёмкости клиента. Если эти вызовы оставить у
медианного поставщика, гибрид на клиентах класса Max дешевле его только в 1.6--2.6 раза. Клиенты класса Ultra обслуживают все 1\,250 вызовов
сами и обходятся в \$37--38, на 4--9\% дороже самого дешёвого API при хранении по политике протокола и в \$53--54
при осторожной политике.

\begin{table}[t]
\centering
\caption{Гибрид на человека в месяц, тяжёлый профиль (2\,500 вызовов в день). Вызовы сверх ёмкости клиента оплачены
по цене самого дешёвого поставщика API. Хранение копий сессий на сервере по политике протокола; осторожная политика
хранения добавляет \$9--17. Ёмкость считана для четырёх потоков генерации за восьмичасовой день; ёмкость M3 Max рассчитана по замеренным кривым, остальные клиенты оценены по вычислительной мощности и полосе
памяти.}
\label{tab:econ}
\footnotesize
\begin{tabular}{lrrrr}
\toprule
Клиент & Ёмкость, & $s = 20\%$, & $s = 50\%$, & $s = 50\%$, \\
 & вызовов в день & Мак уже есть & Мак уже есть & Мак куплен \\
\midrule
M3 Max 128 ГБ & 616 & \$27.3 & \$48.6 & \$96.1 \\
M4 Max 128 ГБ & 833 & \$25.8 & \$46.7 & \$94.1 \\
M5 Max 64 ГБ & 883 & \$24.1 & \$43.6 & \$78.8 \\
M5 Max 128 ГБ & 935 & \$24.1 & \$43.4 & \$100.3 \\
M5 Ultra 96 ГБ & 1\,780 & \$22.6 & \$38.3 & \$112.8 \\
M5 Ultra 256 ГБ & 1\,814 & \$22.0 & \$36.8 & \$183.0 \\
\midrule
API, медианный поставщик & & \$101.7 & \$254.2 & \$254.2 \\
API, самый дешёвый поставщик & & \$14.1 & \$35.2 & \$35.2 \\
\bottomrule
\end{tabular}
\end{table}

\textbf{Доля открытой модели.} Чем больше доля, тем больше гибрид экономит против медианного поставщика и тем раньше
упирается в ёмкость клиента. Пользователю с 800 вызовами в день ёмкости хватает и при доле 50\%. Долю 100\% мы
приводим как верхнюю границу: тяжёлому пользователю, который все вызовы отправляет открытой модели, понадобилось бы примерно три M5 Max, около четырёх M3 Max или около полутора Ultra.

\textbf{Покупка клиента.} Если машину покупают ради гибрида, к цене добавляется доплата за неё сверх обычного
ноутбука. Для M5 Max с 64~ГБ это \$2\,600, которые за четыре года при остаточной стоимости 35\% дают \$35 в месяц.
Тогда гибрид при доле 50\% стоит \$79 против \$35 у самого дешёвого API и \$254 у медианного. Против медианного
поставщика покупка окупается уже при доле 20\%: \$59 против \$102. Рисунок~\ref{fig:payback} показывает, за сколько
месяцев экономия против медианного поставщика возвращает доплату и машину целиком, когда вызовы сверх ёмкости
клиента уходят к самому дешёвому поставщику. Тогда при доле 50\% доплата возвращается за 12 месяцев, Мак целиком
за 24; при доле 20\% за 34 и 66 месяцев. Если вызовы сверх ёмкости оставить у медианного поставщика, при доле 50\%
доплата возвращается за 18 месяцев, Мак целиком за 35. Против самого дешёвого поставщика гибрид дороже,
и доплата не окупается. Доплату обычно платят к рабочему ноутбуку, который
специалисту нужен и так, и кроме экономии на вызовах она покупает память под модель и контекст, которые иначе не
помещаются, полосу памяти для быстрой генерации, несколько агентов одновременно и работу без поставщика. Против
самого дешёвого API доплату может окупить только эта польза: при доле 50\% гибрид с купленным M5 Max дороже его на
\$44 в месяц. Сколько полезной работы даёт эта разница, покажет эксплуатация.

\begin{figure}[t]
\centering
\begin{tikzpicture}
\begin{axis}[
  ybar, bar width=11pt, width=0.62\textwidth, height=50mm, ymin=0, ymax=74,
  symbolic x coords={20,50,100}, xtick=data, xticklabels={{$s = 20\%$},{$s = 50\%$},{$s = 100\%$}},
  ylabel={месяцев}, enlarge x limits=0.25,
  nodes near coords={\pgfmathprintnumber[fixed,precision=0]{\pgfplotspointmeta}},
  nodes near coords style={font=\scriptsize},
  tick label style={font=\footnotesize}, label style={font=\footnotesize},
  legend style={font=\footnotesize,draw=none,at={(0.5,-0.2)},anchor=north,legend columns=2,column sep=6pt}]
\addplot[fill=white,draw=black] table[x=share,y=premium_months,col sep=comma] {data/payback.csv};
\addlegendentry{доплата за M5 Max 64~ГБ, \$2\,600}
\addplot[pattern=north east lines,draw=black] table[x=share,y=whole_months,col sep=comma] {data/payback.csv};
\addlegendentry{машина целиком, \$5\,099}
\draw[dashed] ({rel axis cs:0,0}|-{axis cs:20,48}) -- node[pos=0.5,above,font=\scriptsize]{четыре года} ({rel axis cs:1,0}|-{axis cs:20,48});
\end{axis}
\end{tikzpicture}
\caption{Окупаемость клиента против медианного поставщика API: доплата за M5 Max 64~ГБ сверх обычного ноутбука и
машина целиком, тяжёлый профиль. Вызовы сверх ёмкости клиента уходят к самому дешёвому поставщику; если оставить их
у медианного, при доле 50\% доплата возвращается за 18 месяцев, машина целиком за 35. Месяцы без учёта остаточной
стоимости. Против самого дешёвого поставщика гибрид дороже на \$8--10 в месяц, и покупка не окупается.}
\label{fig:payback}
\end{figure}

\textbf{Команды от 5 до 5000 человек.} Рисунок~\ref{fig:teams} сравнивает гибрид со своим серверным кластером. Узел
B300 покупается целиком, поэтому для малой команды свой кластер очень дорог: на пять человек \$3\,370 на человека в
месяц. С ростом команды цена падает: на 500 человек свой кластер уже дешевле медианного поставщика API при обеих долях
(\$67 против \$102 при доле 20\%, \$169 против \$254 при доле 50\%), на 50 человек ещё дороже. Самого дешёвого
поставщика API свой кластер не догоняет ни при каком размере команды до 5000 человек. Гибрид от размера команды не зависит, потому что серверный prefill здесь покупается долями у общего поставщика по
часовой цене B300 при загрузке 70\%; команде со своим выделенным сервером такая цена недоступна
(таблица~\ref{tab:rentbuy}). Малой команде вместо целого узла B300 подходят
серверы на RTX PRO 6000. Чисто серверный режим на таких серверах мы считаем для закрытого контура, где вся работа идёт на
открытой модели (абзац о закрытом контуре ниже). Команде из пяти человек нужна одна такая карта при лёгкой нагрузке и четыре,
по две в сервере, при тяжёлой, и свой сервер обходится ей в \$145--506 на человека в месяц против \$3\,370 на человека с узлом B300.
Команде из двадцати человек нужна одна карта при лёгкой нагрузке и четырнадцать, по две в сервере, при тяжёлой, и свой
сервер обходится ей в \$36--443 на человека в месяц против \$843 на человека с узлом B300.

\begin{figure}[t]
\centering
\begin{tikzpicture}
\begin{axis}[
  name=left, width=0.49\textwidth, height=58mm, xmode=log, ymode=log, log basis x=10,
  xmin=4, xmax=6500, ymin=10, ymax=5000, xtick={5,10,20,50,500,5000}, xticklabels={5,10,20,50,500,5000},
  xlabel={человек в команде}, ylabel={\$ на человека в месяц}, title={$s = 20\%$},
  tick label style={font=\footnotesize}, label style={font=\footnotesize}, title style={font=\footnotesize},
  legend to name=teamlegend, legend columns=3, legend style={font=\footnotesize,draw=none,column sep=4pt}]
\addplot[black,mark=*,mark size=1.4pt] table[x=team,y=own,col sep=comma] {data/team-cost-s20.csv};
\addlegendentry{свой кластер B300}
\addplot[black,dashed] table[x=team,y=api_median,col sep=comma] {data/team-cost-s20.csv};
\addlegendentry{API, медианный}
\addplot[black,densely dotted,thick] table[x=team,y=api_cheapest,col sep=comma] {data/team-cost-s20.csv};
\addlegendentry{API, самый дешёвый}
\addplot[black,very thick] table[x=team,y=hybrid_owned,col sep=comma] {data/team-cost-s20.csv};
\addlegendentry{гибрид, Мак уже есть}
\addplot[black,dash dot,thick] table[x=team,y=hybrid_bought,col sep=comma] {data/team-cost-s20.csv};
\addlegendentry{гибрид, Мак куплен}
\end{axis}
\begin{axis}[
  at={(left.east)}, anchor=west, xshift=12mm, width=0.49\textwidth, height=58mm, xmode=log, ymode=log,
  xmin=4, xmax=6500, ymin=10, ymax=5000, xtick={5,10,20,50,500,5000}, xticklabels={5,10,20,50,500,5000},
  xlabel={человек в команде}, title={$s = 50\%$},
  tick label style={font=\footnotesize}, label style={font=\footnotesize}, title style={font=\footnotesize}]
\addplot[black,mark=*,mark size=1.4pt] table[x=team,y=own,col sep=comma] {data/team-cost-s50.csv};
\addplot[black,dashed] table[x=team,y=api_median,col sep=comma] {data/team-cost-s50.csv};
\addplot[black,densely dotted,thick] table[x=team,y=api_cheapest,col sep=comma] {data/team-cost-s50.csv};
\addplot[black,very thick] table[x=team,y=hybrid_owned,col sep=comma] {data/team-cost-s50.csv};
\addplot[black,dash dot,thick] table[x=team,y=hybrid_bought,col sep=comma] {data/team-cost-s50.csv};
\end{axis}
\end{tikzpicture}

\vspace{2pt}
\ref*{teamlegend}
\caption{Стоимость на человека в месяц в зависимости от размера команды, тяжёлый профиль, клиент M5 Max 64~ГБ.
Свой кластер собран из целых узлов B300 по восемь GPU и обслуживает вызовы открытой модели целиком на сервере. Гибрид
покупает серверный prefill долями у общего поставщика по часовой цене B300 при загрузке 70\%.}
\label{fig:teams}
\end{figure}

\textbf{Аренда или покупка сервера.} Своё железо нужно гибриду, когда команда держит собственный prefill-сервер,
например ради контроля над данными. Небольшие команды обычно работают удалённо, и сервер им удобнее держать в
колокации. Таблица~\ref{tab:rentbuy} сравнивает почасовую аренду RTX PRO 6000 с покупкой такого же
сервера в колокацию при доле 50\%. При рабочем дне в 10 часов команде из 5--10 человек аренда дешевле покупки при любой из рассмотренных остаточных
стоимостей: даже при 80\% покупка дороже на 7\%. На 50 человек аренда (\$24) почти равна покупке только при
остаточной стоимости 80\% (\$25); при 50 и 20\% покупка дороже в 1.4 и 1.8 раза. Если та же работа распределена на
круглые сутки, например между часовыми поясами, хватает меньшего числа серверов, и покупка выигрывает при любом
размере команды.

\begin{table}[t]
\centering
\caption{Серверная часть гибрида на человека в месяц при доле 50\%, тяжёлый профиль. Аренда RTX PRO 6000 по
\$2.20 в час; покупка сервера с одной RTX PRO 6000 в колокацию, амортизация за четыре года с остаточной стоимостью
20, 50 или 80\%. В строках с 24 ч та же дневная нагрузка, распределённая на сутки в рабочие дни. В скобках число серверов.
Хранение копий сессий по политике протокола входит в обе колонки.}
\label{tab:rentbuy}
\footnotesize
\begin{tabular}{lrrrrr}
\toprule
Команда, нагрузка & Аренда & Покупка, 20\% & Покупка, 50\% & Покупка, 80\% \\
\midrule
5 человек, 10 ч в день & \$98 (1) & \$194 (1) & \$149 (1) & \$104 (1) \\
5 человек, 24 ч & \$231 (1) & \$198 (1) & \$153 (1) & \$108 (1) \\
10 человек, 10 ч в день & \$51 (1) & \$100 (1) & \$77 (1) & \$55 (1) \\
10 человек, 24 ч & \$120 (1) & \$104 (1) & \$81 (1) & \$59 (1) \\
50 человек, 10 ч в день & \$24 (2) & \$43 (2) & \$34 (2) & \$25 (2) \\
50 человек, 24 ч & \$31 (1) & \$28 (1) & \$23 (1) & \$19 (1) \\
\bottomrule
\end{tabular}
\end{table}

\textbf{Закрытый контур.} Часть организаций не может отправлять данные за свой периметр: ни облачным API, ни
стороннему поставщику prefill. Для них доля 100\% становится обычным случаем, и сравнивать приходится три варианта
внутри периметра: одни клиентские машины, свой серверный кластер и гибрид со своими prefill-узлами. Вызовы, которые
клиент не успевает сгенерировать, гибрид отправляет на генерацию своему же серверу. Таблица~\ref{tab:perimeter}
показывает варианты для трёх профилей нагрузки и клиента M5 Max 64~ГБ, который у сотрудника уже есть.

Одни клиентские машины почти ничего не стоят, но пользователю с 800 вызовами в день они не успевают ответить на 5\%
вызовов, а каждый вызов ждёт prefill в среднем 6.5 секунды. Гибрид отвечает на все вызовы, и prefill занимает
у него от 0.8 до 1.9 секунды. Для средней нагрузки гибрид в 2.9--3.9 раза дешевле своего серверного кластера на 50,
500 и 5000 человек, при осторожной политике хранения в 2.4--3.2 раза. Для лёгкой нагрузки на 50 и 5000 человек
гибрид дешевле кластера на 30--39\%, при осторожной политике на 4--22\%. Серверы на RTX PRO 6000 в своей стойке мы
рассматриваем для команд до 100 человек, а узлы B300 начиная с 50, и движок выбирает самый дешёвый из допустимых
вариантов. По этому правилу на 500 человек при лёгкой нагрузке оба варианта получают по одному узлу B300, и свой
кластер обходится в \$34 на человека в месяц против \$46 у гибрида, который сверх узла платит за электричество
клиентов, обслуживание и хранение копий сессий. Вне этого правила, на серверах с двумя RTX PRO 6000, гибрид на 500
человек стоил бы около \$17 (при осторожной политике \$27) против \$35 у своего кластера на тех же серверах и \$34
на узле B300. Экономия немонотонна, потому что серверы и узлы
добавляются целиком. При тяжёлой нагрузке клиент не успевает генерировать, большая часть вызовов уходит на генерацию
на сервере, и гибрид дешевле кластера на 6--26\%, при осторожной политике хранения на 4--24\%. Команде из
пяти человек свой сервер обходится дорого в обоих вариантах: от \$145 на человека в месяц при лёгкой нагрузке.

\begin{table}[t]
\centering
\caption{Закрытый контур: всё внутри периметра, доля открытой модели 100\%, клиент M5 Max 64~ГБ уже есть. Стоимость на
человека в месяц; для одних клиентских машин указана доля вызовов, которую они успевают обслужить за восьмичасовой
день. Сервер: одна-две RTX PRO 6000 в своей стойке для команд до 100 человек или узлы B300 для команд от 50,
остаточная стоимость 20\%; из допустимых вариантов взят самый дешёвый. Гибрид дан с хранением
по политике протокола и по осторожной политике. Место в стойке уже есть, ожидание prefill посчитано без очередей,
парк серверов рассчитан на загрузку не выше 70\%.}
\label{tab:perimeter}
\footnotesize
\begin{tabular}{rrrrrr}
\toprule
Вызовов в день & Команда & Свой кластер & Гибрид & Гибрид, осторожно & Одни клиенты: обслужено \\
\midrule
200 & 50 & \$43 & \$27 & \$34 & 100\% \\
200 & 500 & \$34 & \$46 & \$53 & 100\% \\
200 & 5000 & \$27 & \$19 & \$26 & 100\% \\
800 & 50 & \$145 & \$37 & \$45 & 95\% \\
800 & 500 & \$135 & \$46 & \$53 & 95\% \\
800 & 5000 & \$104 & \$29 & \$43 & 95\% \\
2\,500 & 50 & \$337 & \$316 & \$323 & 31\% \\
2\,500 & 500 & \$337 & \$248 & \$255 & 31\% \\
2\,500 & 5000 & \$320 & \$238 & \$252 & 31\% \\
\bottomrule
\end{tabular}
\end{table}

\textbf{Как это будет меняться.} Облачный пользователь получает прогресс моделей в виде снижения цены: по оценке
Epoch AI стоимость фиксированного уровня способностей с 2023 года падает примерно в 13 раз в год
\cite{epoch2026price}. Пользователь гибрида получает его в виде новых способностей на том же железе: модель класса 27B,
которая помещается на тот же Мак, за полгода выросла по индексу Artificial Analysis с 34.6 до 52.0 \cite{aaindex}
(рисунок~\ref{fig:progress}). Цены API снижаются, и уже сегодня самый дешёвый поставщик обходится дешевле гибрида на всех
рассмотренных клиентах. Открытые модели умнеют, доля открытой модели растёт, и с ней растёт экономия против медианного
поставщика. Контекст становится глубже, и налог на поддержание кэша растёт вместе с глубиной, а работа гибридного
сервера растёт медленнее: он заново читает только новые токены и ответы клиента, хотя цена каждого нового токена тоже
растёт с глубиной. Поэтому главный рычаг гибрида находится на глубоком контексте
(раздел~\ref{sec:million}).

\begin{figure}[t]
\centering
\begin{tikzpicture}
\begin{axis}[
  name=price, width=0.47\textwidth, height=48mm, ymode=log, xmin=0, xmax=8.5, ymin=0.1, ymax=1.2,
  xtick={0,2,4,6,8}, xticklabels={фев,апр,июн,авг,окт}, ytick={0.1,0.2,0.5,1}, yticklabels={0.1,0.2,0.5,1},
  ylabel={цена, доля от февраля}, title={API: цена фиксированного уровня},
  tick label style={font=\footnotesize}, label style={font=\footnotesize}, title style={font=\footnotesize}]
\addplot[black,dashed,thick,domain=0:8,samples=40] {13^(-x/12)};
\end{axis}
\begin{axis}[
  at={(price.east)}, anchor=west, xshift=14mm, width=0.47\textwidth, height=48mm, xmin=0, xmax=8.5, ymin=30, ymax=56,
  xtick={0,2,4,6,8}, xticklabels={фев,апр,июн,авг,окт}, ylabel={индекс AA v4.1.1},
  title={Гибрид: модель 27B на том же Маке},
  tick label style={font=\footnotesize}, label style={font=\footnotesize}, title style={font=\footnotesize}]
\addplot[black,thick,mark=*,mark size=1.6pt] coordinates {(0.75,34.6) (2.7,37.7) (6.45,52.0)};
\node[font=\scriptsize,anchor=north west] at (axis cs:0.75,34.6) {Qwen3.5};
\node[font=\scriptsize,anchor=south east] at (axis cs:2.7,37.7) {Qwen3.6};
\node[font=\scriptsize,anchor=south east] at (axis cs:6.45,52.0) {Qwen3.8};
\end{axis}
\end{tikzpicture}
\caption{Как пользователь получает прогресс моделей в 2026 году. Слева: цена фиксированного уровня способностей через
API при снижении в 13 раз в год \cite{epoch2026price}, нормирована на февраль. Справа: индекс Artificial Analysis v4.1.1
моделей класса 27B, которые помещаются в одну и ту же машину с 32~ГБ памяти \cite{aaindex}; значение для Qwen3.5-27B
оценочное.}
\label{fig:progress}
\end{figure}

\textbf{Вывод.} Против медианного поставщика API гибрид на уже купленной машине класса Max, которая сама обслуживает
почти все вызовы открытой модели (доля 20\%), дешевле в 2.8--4.2 раза при любой из рассмотренных политик хранения.
Клиенты класса Ultra при доле 50\% обслуживают все вызовы сами и дешевле медианного поставщика в 4.7--6.9 раза в
зависимости от политики хранения. При доле 50\% машина класса Max обслуживает от 49 до 75\% вызовов, и соотношение в 5.2--5.9 раза
включает переход остальных вызовов к самому дешёвому поставщику. Свой кластер из узлов B300 догоняет медианного
поставщика только на сотнях человек. Самый дешёвый поставщик API дешевле гибрида на всех рассмотренных клиентах и
при обеих долях. В закрытом контуре стороннего API нет, и там гибрид при средней нагрузке в 2.4--3.9 раза дешевле
своего кластера, при тяжёлой на 6--26\% (при осторожной политике хранения на 4--24\%), а при лёгкой дешевле на 50
и 5000 человек и дороже на 500: серверы на RTX PRO 6000 мы рассматриваем только для команд до 100 человек, и на 500
человек оба варианта получают по одному узлу B300. Покупать Мак ради гибрида имеет смысл против медианного поставщика:
при доле 50\% доплата возвращается за год, если вызовы сверх ёмкости клиента уходят к самому дешёвому поставщику,
и за полтора года, если к медианному.
